\section{Simulation Protocols and Metrics}
\label{app:protocols}

This appendix summarizes simulation conditions, replication counts, schedules, and metrics.

\subsection{Simulation conditions}

All reported experiments use one focal active-inference agent with four partners ($K=4$), per-partner generative models, graded investment actions, and environment-side partner policies (Appendix~\ref{app:pomdp}). Conditions differ by assignment mode, scheduled partner changes, and precision/profile manipulations within the graded payoff regime.

\textbf{Open graded decision regime.} The graded regime uses a discretized investment level while keeping the same latent type--stance structure. This is the primary deployment-contrast setting (Section~\ref{sec:deployment}): affect, no-affect, and tracked-only variants are compared on policy entropy and cumulative payoff.

\textbf{Predictability-over-value locality probe.} A 100-round partner-choice graded environment with partner-local, shared-$\beta$, and no-affect variants tests whether $\beta_k$-derived policy precision tracks partner-response likelihood surprisal more strongly than realized payoff (Section~\ref{sec:predictability_value}). The payoff panel reports the same active-encounter analysis window used for the locality-probe correlation readouts.

\textbf{Partner-choice regime.} The agent first selects a partner, then selects an investment level toward that partner each round. Partner selection is where engagement reorganization and social-allocation readouts are defined (Section~\ref{sec:partner_choice}).

\textbf{Abrupt betrayal.} A partner-choice graded environment with scheduled stance betrayal on a previously cooperative partner tests post-switch reallocation, entropy, payoff, and joint accuracy (Section~\ref{sec:betrayal}).

\textbf{Precision-dynamics/profile variation.} Computational profile variants vary $\beta_k$ amplitude and stability by manipulating precision gain $\alpha$ and prior model fitness in open graded, betrayal, and partner-choice environments (Section~\ref{sec:phenotypes}; Appendix~\ref{app:extended_results}).

\textbf{Forgiveness.} The trust-repair protocol extends abrupt betrayal by allowing the betrayed partner to return to a cooperative/trusting state, testing reengagement and confidence recovery after restored reliability (Section~\ref{sec:phenotypes}).

\subsection{Seed counts and episode lengths}

Unless otherwise stated in the main text, episode length is 200 rounds. The abrupt-betrayal condition uses 120 rounds because post-switch readouts dominate interpretation. Profile-scale tables use 20 seeds per condition. Central behavioral comparisons use 30 seeds per condition. Exploratory diagnostic runs used three seeds per variant and were used only for implementation validation; all main-text claims are based on the confirmation-scale or profile-scale runs reported in the relevant sections.

Representative configured scales:

\begin{table}[H]
\centering
\caption{Representative episode lengths and replication counts by experiment family.}
\label{tab:protocol_scales}
\begin{tabular}{llrr}
\toprule
Experiment family & Setting & Rounds & Seeds \\
\midrule
Open deployment & graded, random/open choice & 200 & 30 \\
Predictability-over-value locality probe & graded partner choice & 100 & 30 \\
Partner allocation & partner choice & 200 & 30 \\
Abrupt betrayal & scheduled partner-choice betrayal & 120 & 30 \\
Precision perturbation & graded perturbations & 200 & 30 \\
Profile program & sweeps, factorial, forgiveness & 200 & 20 \\
\bottomrule
\end{tabular}
\end{table}

\paragraph{Supplementary materials.}
The supplementary repository contains the simulation environments, affect-runtime variants, experiment configuration files, analysis scripts, and figure-generation code used for the reported experiments.

\subsection{Betrayal schedule}

The canonical abrupt-betrayal protocol (\texttt{betrayal\_choice}) uses graded payoffs, agent-chosen partners, four partners with fixed initial types and stances, and no stochastic type switching ($p_{\mathrm{switch}}=0$). Partner P0 begins as a cooperator in a trusting stance. At round 31, P0 undergoes a scheduled stance switch to hostile while other partners retain their initial configurations. Primary readouts are post-switch cumulative payoff, policy entropy, partner reallocation, and joint partner-type accuracy over the remaining episode.

The profile-scale betrayal environments used in the alpha sweep and prior-gain factorial apply the same single-switch structure within 200-round episodes for recovery-time and quadrant metrics. The forgiveness profile extends this schedule: rounds 1--80 cooperative baseline, rounds 81--120 betrayal on P0, rounds 121--200 reversion to cooperative/trusting on P0.

\subsection{Partner-choice allocation metrics}

Partner-choice readouts quantify how precision reshapes \emph{who} the agent engages, not only how it acts toward a fixed partner.

\textbf{Selection share.} Per-partner selection rate is the fraction of rounds on which each partner is chosen. Partner-choice reallocation is summarized qualitatively in Section~\ref{sec:partner_choice}.

\textbf{Selection concentration.} The Gini coefficient of the partner-selection distribution summarizes how concentrated engagement is across the social context (used in the profile alpha sweep and prior-gain factorial).

\textbf{Post-betrayal allocation.} Post-switch windows report reallocation away from the switched partner and toward alternatives, including post-betrayal P0 selection and high-investment rates in profile analyses.

\textbf{Policy entropy under choice.} Mean $q_\pi$ entropy in partner-choice regimes captures how sharply the agent commits to partner--action policies (reported comparatively for affect versus no-affect controls).

\subsection{Entropy, payoff, and accuracy readouts}

\textbf{Total payoff.} Cumulative focal-agent payoff over the episode (or over defined post-switch windows for betrayal analyses).

\textbf{Policy entropy.} Mean $q_\pi$ entropy (`q\_pi\_entropy` in logged outputs) measures policy-commitment sharpness; lower entropy indicates sharper policy selection under precision modulation.

\textbf{Joint partner-type accuracy.} Mean joint correctness of inferred partner type (`mean\_joint\_accuracy`) summarizes epistemic state inference; it is reported separately from payoff because precision modulation can change deployment and accuracy independently of cumulative reward.

\textbf{$\beta_k$ range.} Unless otherwise specified, $\beta_k$ range denotes the mean within-episode range of the posterior mean $\mathbb{E}_{q}[\beta_k]$, averaged across relevant partners and seeds. It is an observed precision-dynamics amplitude metric, not the fixed support of the categorical $\beta_k$ variable.

\textbf{Betrayal timing metrics.} Recovery time and related latency readouts measure rounds after a switch until selection or payoff returns toward baseline thresholds (profile alpha-sweep, prior-gain factorial, and betrayal analyses).

\subsection{Replication summary}

Central behavioral comparisons in Section~\ref{sec:results} use 30 seeds per condition. Profile-scale analyses use 20 seeds per condition.
